% ============================================================
%  Pin Mapping: Raspberry Pi Zero 2 W  <-->  Adafruit 1651
%  2.8" TFT Touch Shield (ILI9341 + TSC2007)
%  Custom PCB Design Reference  --  DIY Kit
% ============================================================
\documentclass[11pt,a4paper]{article}

\usepackage[a4paper, margin=2cm]{geometry}
\usepackage{booktabs}
\usepackage{tabularx}
\usepackage{longtable}
\usepackage{array}
\usepackage{xcolor}
\usepackage{colortbl}
\usepackage{hyperref}
\usepackage{parskip}
\usepackage{microtype}
\usepackage{fancyhdr}
\usepackage{titlesec}
\usepackage{mdframed}
\usepackage{enumitem}
\usepackage{multirow}
\usepackage{amsmath}

% ── Colours ──────────────────────────────────────────────────
\definecolor{hdr}{RGB}{30,70,120}
\definecolor{rowA}{RGB}{235,242,252}
\definecolor{rowB}{RGB}{248,251,255}
\definecolor{warn}{RGB}{255,243,205}
\definecolor{warnborder}{RGB}{200,150,0}
\definecolor{notebg}{RGB}{237,247,237}
\definecolor{noteborder}{RGB}{60,140,60}
\definecolor{oddcol}{RGB}{220,235,255}
\definecolor{evencol}{RGB}{220,250,225}
\definecolor{pwrcol}{RGB}{255,230,220}

% ── Column types ─────────────────────────────────────────────
\newcolumntype{L}[1]{>{\raggedright\arraybackslash}p{#1}}
\newcolumntype{C}[1]{>{\centering\arraybackslash}p{#1}}

% ── Section styling ───────────────────────────────────────────
\titleformat{\section}{\large\bfseries\color{hdr}}{}{0em}{}[\titlerule]
\titleformat{\subsection}{\normalsize\bfseries\color{hdr}}{}{0em}{}

% ── Header / footer ───────────────────────────────────────────
\pagestyle{fancy}
\fancyhf{}
\lhead{\small RPi Zero 2W $\leftrightarrow$ Adafruit 1651 Pin Mapping}
\rhead{\small DIY Kit PCB Reference}
\cfoot{\small Page \thepage}
\renewcommand{\headrulewidth}{0.4pt}

% ── Note / warning boxes ──────────────────────────────────────
\newmdenv[
  backgroundcolor=notebg,
  linecolor=noteborder,
  linewidth=1.2pt,
  leftmargin=0pt, rightmargin=0pt,
  innerleftmargin=8pt, innerrightmargin=8pt,
  innertopmargin=6pt, innerbottommargin=6pt,
  skipabove=6pt, skipbelow=6pt
]{notebox}

\newmdenv[
  backgroundcolor=warn,
  linecolor=warnborder,
  linewidth=1.2pt,
  leftmargin=0pt, rightmargin=0pt,
  innerleftmargin=8pt, innerrightmargin=8pt,
  innertopmargin=6pt, innerbottommargin=6pt,
  skipabove=6pt, skipbelow=6pt
]{warnbox}

% ── Helpers ───────────────────────────────────────────────────
\newcommand{\fixed}{\textbf{FIXED}}
\newcommand{\soft}{\textit{soft}}
\newcommand{\na}{---}

% =============================================================
\begin{document}

\begin{center}
  {\LARGE\bfseries\color{hdr}
    Pin Mapping Reference\\[4pt]
    Raspberry Pi Zero 2~W $\leftrightarrow$ Adafruit 1651}\\[6pt]
  {\large 2.8\,'' TFT Touch Shield with Resistive Touch (TSC2007)}\\[4pt]
  {\normalsize Custom PCB Design --- DIY Kit\\
   \textit{Point-to-point wiring between RPi header and 1651 shield header}}\\[2pt]
  {\small\color{gray}
    Sources: Adafruit Learning System (learn.adafruit.com/adafruit-2-8-tft-touch-shield-v2),\\
    Raspberry Pi GPIO documentation
    (\href{https://www.raspberrypi.com/documentation/computers/raspberry-pi.html\#gpio}%
    {raspberrypi.com/documentation/computers/raspberry-pi.html\#gpio})}
\end{center}

\vspace{4pt}\hrule\vspace{10pt}

% ── RPi GPIO Quick Reference ──────────────────────────────────
\section{RPi Zero 2~W GPIO Peripheral Reference}

The following functions are fixed by the BCM2710A1 SoC and cannot be
remapped in software.  All GPIO pins operate at \textbf{3.3\,V logic}.
Applying 5\,V to any GPIO pin will permanently damage the board.

\begin{notebox}
\textbf{Source:}
\href{https://www.raspberrypi.com/documentation/computers/raspberry-pi.html\#gpio}%
{raspberrypi.com --- GPIO documentation}

\smallskip
\textbf{PWM (pulse-width modulation)}
\begin{itemize}[noitemsep,topsep=2pt]
  \item Software PWM available on all pins
  \item Hardware PWM available on \textbf{GPIO12, GPIO13, GPIO18, GPIO19}
\end{itemize}

\textbf{SPI}
\begin{itemize}[noitemsep,topsep=2pt]
  \item SPI0: MOSI (GPIO10); MISO (GPIO9); SCLK (GPIO11);
        CE0 (GPIO8); CE1 (GPIO7)
  \item SPI1: MOSI (GPIO20); MISO (GPIO19); SCLK (GPIO21);
        CE0 (GPIO18); CE1 (GPIO17); CE2 (GPIO16)
\end{itemize}

\textbf{I2C}
\begin{itemize}[noitemsep,topsep=2pt]
  \item Data: GPIO2 \quad Clock: GPIO3
  \item EEPROM Data: GPIO0 \quad EEPROM Clock: GPIO1
        \quad \textit{(do not use for general I/O)}
\end{itemize}

\textbf{Serial (UART)}
\begin{itemize}[noitemsep,topsep=2pt]
  \item TX: GPIO14 (physical pin~8) \quad RX: GPIO15 (physical pin~10)
  \item \textit{These are reserved for the Linux serial console by default.
        Avoid using them for general I/O in a DIY kit.}
\end{itemize}
\end{notebox}

% ── Power Pin Note ────────────────────────────────────────────
\subsection{Power and Ground Pin Notes}

\begin{warnbox}
\textbf{5V pins --- Pin~2 and Pin~4 (both even side):}
Both pins are electrically identical --- connected directly to the
incoming micro-USB supply rail with no intervening regulation or
current limiting beyond the board's polyfuse.
The duplication is for PCB routing convenience only.
Either may be used; for a DIY kit, Pin~2 is conventional.
Current capacity is limited only by your power supply and the board's
polyfuse (typically rated around 2\,A for the Zero 2\,W).

\smallskip
\textbf{3.3V pins --- Pin~1 (odd) and Pin~17 (odd):}
Both are outputs of the same onboard 3.3\,V regulator.
They are electrically identical.
The \textbf{total} current available from the 3.3\,V rail
(shared with all GPIO output drivers and any attached peripherals)
is approximately \textbf{50\,mA}.
The 1651 shield has its own onboard 3.3\,V regulator fed from the 5\,V
rail, so it does \emph{not} draw from the RPi's 3.3\,V supply.
Do not power the 1651 from the RPi 3.3\,V pin.

\smallskip
\textbf{Ground pins:}
Physical pins 6, 9, 14, 20, 25, 30, 34, and 39 are all ground.
They are all identical. Use the ground pin physically closest to
the signal pins being routed for best noise performance.
For a two-connector PCB layout (odd-side and even-side headers),
use Pin~9 or Pin~25 for the odd-side (1651) ground return,
and Pin~6, 14, or 20 for the even-side ground return.
\end{warnbox}

% =============================================================
\section{RPi Zero 2~W --- Complete Pin Allocation Table}

\textbf{Legend:} \colorbox{oddcol}{\strut Odd-side pin} \quad
\colorbox{evencol}{\strut Even-side pin} \quad
\colorbox{pwrcol}{\strut Power / Ground}

\smallskip
Column headers: \textbf{Phys} = physical header pin number;
\textbf{BCM} = Broadcom GPIO number used in Python code;
\textbf{Dir} = I=input, O=output, PWM, PWR;
\textbf{Assignment} = function in this design.

\vspace{6pt}

\setlength{\tabcolsep}{4pt}
\renewcommand{\arraystretch}{1.30}

\begin{longtable}{C{0.9cm} C{1.3cm} C{0.8cm} L{3.5cm} L{7.5cm}}
  \toprule
  \rowcolor{hdr}
  \color{white}\textbf{Phys} &
  \color{white}\textbf{BCM} &
  \color{white}\textbf{Dir} &
  \color{white}\textbf{Assignment} &
  \color{white}\textbf{Notes} \\
  \midrule
  \endfirsthead
  \toprule
  \rowcolor{hdr}
  \color{white}\textbf{Phys} &
  \color{white}\textbf{BCM} &
  \color{white}\textbf{Dir} &
  \color{white}\textbf{Assignment} &
  \color{white}\textbf{Notes} \\
  \midrule
  \endhead
  \bottomrule
  \endfoot

  %% ── ODD SIDE ─────────────────────────────────────────────
  \multicolumn{5}{l}{\cellcolor{oddcol}\textbf{Odd-side pins --- 1651 Display and Touch (point-to-point wiring to shield)}} \\
  \midrule

  \rowcolor{pwrcol}
  1 & \na & PWR & 3.3\,V & \textit{Not used for 1651. Available for other 3.3\,V loads
    if total rail current stays under 50\,mA.} \\

  \rowcolor{oddcol}
  3 & GPIO2 & I/O & I2C SDA $\to$ TSC2007 &
    \fixed\ by SoC. I2C bus 1 data.
    Enable with \texttt{dtparam=i2c\_arm=on} in \texttt{/boot/config.txt}.
    The 1651 shield has onboard 4.7\,k$\Omega$ pull-ups on SDA/SCL --- do not add
    additional pull-ups on the PCB. See protective resistor note (Section~\ref{sec:passive}). 
    [To J1-Pin2]\\

  \rowcolor{rowB}
  5 & GPIO3 & I/O & I2C SCL $\to$ TSC2007 &
    \fixed\ by SoC. I2C bus 1 clock.
    Same pull-up note as SDA above. 
    [To J1-Pin4]\\

  \rowcolor{oddcol}
  7 & GPIO4 & I & TSIRQ (touch interrupt) &
    \soft\ assignment.
    Active-low interrupt from TSC2007.
    Enable internal pull-up in software: \texttt{GPIO.setup(4, GPIO.IN, pull\_up\_down=GPIO.PUD\_UP)}.
    Add 33\,$\Omega$ series resistor on PCB for ESD protection. 
    [To J1-Pin10]\\

  \rowcolor{pwrcol}
  9 & \na & GND & Ground (1651 return) &
    Connect to ground plane. \\

  \rowcolor{oddcol}
  11 & GPIO17 & \na & \textit{Spare} &
  Unallocated. Reserve for future use. 
  Add 33\,$\Omega$ series resistor on PCB. 
  [To J1-Pin14]\\

  \rowcolor{pwrcol}
  17 & \na & PWR & 3.3\,V &
    \textit{Not used for 1651. See power notes above.} \\

  \rowcolor{oddcol}
  19 & GPIO10 & O & SPI0 MOSI $\to$ ILI9341 &
    \fixed\ by SoC (SPI0 MOSI). High-speed data to display.
    Add 33\,$\Omega$ series resistor on PCB.
    [To J1-Pin 13] \\

  \rowcolor{rowB}
  21 & GPIO9 & I & SPI0 MISO $\leftarrow$ ILI9341 &
    \fixed\ by SoC (SPI0 MISO). Read-back from display/SD.
    Add 33\,$\Omega$ series resistor on PCB. 
    [To J1-Pin 14]\\

  \rowcolor{oddcol}
  23 & GPIO11 & O & SPI0 SCLK $\to$ ILI9341 &
    \fixed\ by SoC (SPI0 clock).
    Add 33\,$\Omega$ series resistor on PCB. 
    [To J1-Pin 11]\\

  \rowcolor{pwrcol}
  25 & \na & GND & Ground &
    Connect to ground plane. \\

  \rowcolor{rowB}
  27 & GPIO0 & \na & \textit{PIN REMOVED, DO NOT USE} &
    Reserved for HAT ID EEPROM. Leave unconnected. \\

  \rowcolor{pwrcol}
  29 & GPIO5 & \na & Unassigned &
  May be used as active-low SD Card CS on 1651.
  Add 33\,$\Omega$ series resistor on PCB. 
  [To J1-Pin 8]\\

  \rowcolor{rowB}
  31 & GPIO6 & O & TFT DC (Data/Command) &
  \soft\ assignment. Data/Command select for ILI9341.
  Add 33\,$\Omega$ series resistor on PCB. 
  [To J1-Pin 7]\\

  \rowcolor{oddcol}
  33 & GPIO13 & PWM & Backlight PWM $\to$ 1651 LITE &
    \soft\ assignment. \textbf{Hardware PWM channel 1} (BCM2710 PWM1).
    Connect to 1651 LITE pin (solder the LITE/\#3 jumper on shield).
    Add 33\,$\Omega$ series resistor on PCB.
    Software: use \texttt{RPi.GPIO} PWM or \texttt{pigpio} for precise frequency.  
    [To J1-Pin 9] \\

  \rowcolor{rowB}
  35 & GPIO19 & ? & SPI1 MISO &
    \soft\ assignment. (SPI1 MOSI pin is used for the Shift Register).
    Pin 35 is exposed on J4-Pin 4. It can be used as general IO, or as SPI1 MISO.
    Add 33\,$\Omega$ series resistor on PCB. [To J4-Pin4] \\

  \rowcolor{pwrcol}
  37 & GPIO26 & \na & \textit{Spare} &
  Unallocated. Reserve for future use. [To J1-Pin 16]\\

  \rowcolor{pwrcol}
  39 & \na & GND & Ground &
  Connect to ground plane. \\

  \midrule
  %% ── EVEN SIDE ─────────────────────────────────────────────
  \multicolumn{5}{l}{\cellcolor{evencol}\textbf{Even-side pins --- General I/O (LEDs, Inputs, 1-Wire, MOSFET control)}} \\
  \midrule

  \rowcolor{pwrcol}
  2 & \na & PWR & 5\,V $\to$ 1651 VIN &
    \textbf{Primary 5\,V supply for the 1651 shield.}
    The shield's onboard regulator converts this to 3.3\,V for the ILI9341 and TSC2007.
    Do not use the RPi 3.3\,V rail to power the 1651. 
    [To J1-Pin 1]\\

  \rowcolor{pwrcol}
  4 & \na & PWR & 5\,V $\to$ LED's &
    Electrically identical to Pin~2. Available for other 5\,V loads on the PCB
    (e.g., MOSFET gate driver supply if needed).  
    [To LED \\

  \rowcolor{pwrcol}
  6 & \na & GND & Ground (even-side return) &
   Connect to ground plane. \\

  \rowcolor{evencol}
  8 & GPIO14 & \na & \textit{UART TX --- DO NOT USE} &
    Linux serial console TX by default. Do not use for general I/O in a kit.
    Leave unconnected at header. \\

  \rowcolor{rowB}
  10 & GPIO15 & \na & \textit{UART RX --- DO NOT USE} &
    Linux serial console RX by default. Leave unconnected at header. \\

  \rowcolor{evencol}
  12 & GPIO18 & O & Indicator LED 3 &
    \soft\ assignment. Hardware PWM capable (PWM0) if dimming desired.
    Option Interface header J3: via jumper JP2 (closed) $\to$ J3-Pin4.
    Option SPI1 CE0/Interface header J4: via jumper JP1 (open) $\to$ J4-Pin6  
    Add 470\,$\Omega$ current-limit resistor.\\

  \rowcolor{pwrcol}
  14 & \na & GND & Ground &
    Connect to ground plane.\\

  \rowcolor{evencol}
  16 & GPIO23 & O & Indicator LED 2 &
    \soft\ assignment. General GPIO output.
    Interface header J3 $\to$ J3-Pin4.
    Add 33\,$\Omega$ current-limit resistor in series with LED on PCB. \\

  \rowcolor{rowB}
  18 & GPIO24 & O?I? & 74HC595 RCLK &
    Connect to pin 12 on 74HC595 
    Add 33\,$\Omega$ series ESD resistor on PCB.\\

  \rowcolor{pwrcol}
  20 & \na & GND & Ground &
    Connect to ground plane.\\

  \rowcolor{evencol}
  22 & GPIO25 & O & Control LED 10 (MOSFET gate) &
    500\,$\Omega$ gate resistor and 10K\,$\Omega$ bleed resistor.\\
  
  \rowcolor{evencol}
  24 & GPIO8 & O & TFT CS $\to$ ILI9341 (SPI0 CE0) &
    \fixed\ by SoC (SPI0 CE0). \textbf{Only even-pin forced exception for 1651.}
    Hardware-managed chip select.
    Add 33\,$\Omega$ series resistor on PCB. \\

  \rowcolor{rowB}
  26 & GPIO7 & O & Control LED 9 (MOSFET gate) &
  500\,$\Omega$ gate resistor and 10K\,$\Omega$ bleed resistor.\\
  

  \rowcolor{pwrcol}
  28 & GPIO1 & \na & \textit{PIN REMOVED, DO NOT USE} &
    Reserved for HAT ID EEPROM. Leave unconnected. \\

  \rowcolor{pwrcol}
  30 & \na & GND & Ground &
    Connect to ground plane. \\

  \rowcolor{evencol}
  32 & GPIO12 & PWM & Buzzer &
    Hardware PWM  (PWM0 channel A) 
    Add 1K\,$\Omega$ current-limit resistor.\\

  \rowcolor{pwrcol}
  34 & \na & GND & Ground &
    Additional even-side ground. \\

  \rowcolor{evencol}
  36 & GPIO16 & I/O & 1-Wire bus &
    This is not the default 1-Wire pin.
	\textbf{Requires:} \texttt{dtoverlay=w1-gpio,gpiopin=7} in \texttt{/boot/config.txt}.
	\textbf{Requires 4.7\,k$\Omega$ pull-up to 3.3\,V on PCB}.
	(This is a PCB passive component, not a GPIO internal pull-up --- 
	the 1-Wire protocol requires a strong external pull-up.)\\
		  
  \rowcolor{rowB}
  38 & GPIO20 & O & SPI1-MOSI &
    Connect to Pin 14 of 74HC595.  Also connects to J4-Pin3, where it can be used as SPI1-MOSI or potentially some other output signal
    Add 470\,$\Omega$ gate resistor on PCB. \\

  \rowcolor{pwrcol}
  39 & \na & GND & Ground &
    See odd-side row above. \\

  \rowcolor{evencol}
  40 & GPIO21 & ? & Connect to 74HC595 SRCLCK Pin 11. \\

\end{longtable}

\begin{warnbox}
	These are as-built.
\end{warnbox}

% =============================================================
\section{Adafruit 1651 Shield Pin Mapping Table}
\label{sec:1651}

Point-to-point wiring from each shield pin to the RPi header.
Column 1 is the Arduino/Metro silkscreen label on the 1651 shield.
Passive components (series resistors) sit on the PCB between the RPi
header and the wire going to the 1651 shield pin.

\vspace{6pt}
\renewcommand{\arraystretch}{1.30}

\begin{longtable}{C{1.5cm} C{1.4cm} C{1.4cm} L{2.8cm} L{6.5cm}}
  \toprule
  \rowcolor{hdr}
  \color{white}\textbf{1651 Pin} \newline \color{white}\footnotesize(Arduino label) &
  \color{white}\textbf{RPi Odd} \newline \color{white}\footnotesize(physical/BCM) &
  \color{white}\textbf{RPi Even} \newline \color{white}\footnotesize(physical/BCM) &
  \color{white}\textbf{Description} &
  \color{white}\textbf{Notes} \\
  \midrule
  \endfirsthead
  \toprule
  \rowcolor{hdr}
  \color{white}\textbf{1651 Pin} \newline \color{white}\footnotesize(Arduino label) &
  \color{white}\textbf{RPi Odd} \newline \color{white}\footnotesize(physical/BCM) &
  \color{white}\textbf{RPi Even} \newline \color{white}\footnotesize(physical/BCM) &
  \color{white}\textbf{Description} &
  \color{white}\textbf{Notes} \\
  \midrule
  \endhead
  \bottomrule
  \endfoot

  %% Power
  \rowcolor{pwrcol}
  \textbf{5V / VIN} &
  \na &
  Pin~2 &
  Board power in &
  5\,V from RPi micro-USB rail. Feeds onboard 3.3\,V regulator.
  Do \emph{not} use RPi 3.3\,V.
  No series resistor --- power line. Use adequate trace width on PCB. \\

  \rowcolor{pwrcol}
  \textbf{GND} &
  Pin~9 or 25 &
  Pin~6, 14, or 20 &
  Ground &
  Common ground. Use the ground pin nearest the signal traces being routed. \\

  %% SPI bus
  \rowcolor{oddcol}
  \textbf{ICSP MOSI} &
  Pin~19 / GPIO10 &
  \na &
  SPI0 MOSI $\to$ ILI9341 &
  \fixed. 33\,$\Omega$ series resistor on PCB (RPi side of resistor connects to pin~19;
  shield side connects to ICSP MOSI pin). \\

  \rowcolor{rowB}
  \textbf{ICSP MISO} &
  Pin~21 / GPIO9 &
  \na &
  SPI0 MISO $\leftarrow$ ILI9341 &
  \fixed. 33\,$\Omega$ series resistor on PCB. \\

  \rowcolor{oddcol}
  \textbf{ICSP SCK} &
  Pin~23 / GPIO11 &
  \na &
  SPI0 clock &
  \fixed. 33\,$\Omega$ series resistor on PCB. \\

  \rowcolor{evencol}
  \textbf{D10} --- TFT CS &
  \na &
  Pin~24 / GPIO8 &
  Display chip select (CE0) &
  \fixed\ (SPI0 CE0). Only even-side pin forced for 1651.
  33\,$\Omega$ series resistor on PCB. Active-low. \\

  \rowcolor{oddcol}
  \textbf{D9} --- TFT DC &
  Pin~35 / GPIO19 &
  \na &
  Data/Command select &
  \soft. 33\,$\Omega$ series resistor on PCB. \\

  %% SD card
  \rowcolor{oddcol}
  \textbf{D4} --- SD CS &
  Pin~11 / GPIO17 &
  \na &
  microSD chip select &
  \soft. Active-low. 33\,$\Omega$ series resistor on PCB.
  Only wire if microSD is used; otherwise leave open. \\

  %% Touch / I2C
  \rowcolor{oddcol}
  \textbf{SDA} &
  Pin~3 / GPIO2 &
  \na &
  I2C SDA $\to$ TSC2007 &
  \fixed. The 1651 shield has onboard 4.7\,k$\Omega$ pull-up.
  Do \emph{not} add a second pull-up on PCB.
  33\,$\Omega$ series resistor on PCB for ESD. \\

  \rowcolor{rowB}
  \textbf{SCL} &
  Pin~5 / GPIO3 &
  \na &
  I2C clock $\to$ TSC2007 &
  \fixed. Same pull-up and ESD notes as SDA. \\

  \rowcolor{oddcol}
  \textbf{D2} --- TSIRQ &
  Pin~7 / GPIO4 &
  \na &
  Touch interrupt (active-low) &
  \soft. Solder the TSIRQ/\#2 jumper on 1651 to enable.
  Enable RPi internal pull-up in software.
  33\,$\Omega$ series resistor on PCB. \\

  %% Backlight
  \rowcolor{oddcol}
  \textbf{D3} --- LITE &
  Pin~33 / GPIO13 &
  \na &
  Backlight PWM &
  \soft. Solder the LITE/\#3 jumper on 1651 to enable.
  Hardware PWM channel 1.
  33\,$\Omega$ series resistor on PCB. \\

  %% VddIO note
  \rowcolor{warn}
  \textbf{VddIO jumper} &
  \multicolumn{4}{L{12cm}}{%
    \textbf{IMPORTANT:} The 1651 VddIO jumper must be set to \textbf{3V}
    (solder left pad to centre pad) because the RPi Zero 2\,W is a 3.3\,V logic device.
    This selects the logic reference voltage for the 74LVX245 level shifters on the shield.
    Incorrect setting risks damaged GPIO pins.} \\

\end{longtable}

new version of table

\begin{longtable}{C{1.5cm} C{1.0cm} C{1.0cm} C{1.4cm} L{2.6cm} L{6.9cm}}
	\toprule
	\rowcolor{hdr}
	\color{white}\textbf{1651 GPIO} \newline \color{white}\footnotesize(Arduino label) &
	\color{white}\textbf{1651 Pin} \newline \color{white}\footnotesize(Arduino label) &
	\color{white}\textbf{J1 Pin} \newline \color{white}\footnotesize(header KiCad) &
	\color{white}\textbf{RPi Pin} \newline \color{white}\footnotesize(rpi header) &
	\color{white}\textbf{RPi}\newline\textbf{GPIO} \newline \color{white}\footnotesize(rpi header) &
	\color{white}\textbf{Description}  \\
	\midrule
	\endfirsthead
	\bottomrule
	\endfoot
	
	%% Power
	\rowcolor{pwrcol}
	\textbf{5V / VIN} & (10) & J~1 & pin~2 & V5+ &
	5\,V from RPi micro-USB rail. Feeds onboard 3.3\,V regulator.
	Do \emph{not} use RPi 3.3\,V.
	No series resistor --- power line. Use adequate trace width on PCB. \\
	
	\rowcolor{pwrcol}
	\textbf{GND} & (8,9) & J~3 J~5 & GND & Ground &
	Connects to ground plane on boards. \\
	
	%% SPI reset
	\rowcolor{evencol}
	\textbf{RESET} & RESET & J~14 & pin~11 & GPIO17 &
	\soft. 33\,$\Omega$ series resistor on PCB. 
	Used for hardware reset of ILI9341.
	More reliable than the ILI9341's software reset (0x01).\\
	
	%% SPI bus
	\rowcolor{oddcol}
	\textbf{ICSP MOSI} & D11 & J~13 & pin~19 & GPIO10/SPI0 MOSI &
	\fixed. 33\,$\Omega$ series resistor on PCB at RPi pin 19.
	$\to$ TFT ILI9341 MOSI \\
	  
	%% SPI bus
	\rowcolor{evencol}
	\textbf{ICSP MISO} & D12 & J~12 & pin-21 & GPIO9/SPI0 MISO &
	\fixed. 33\,$\Omega$ series resistor on PCB. 
	Shared SPI0 MISO return path from TFT ILI9341 MISO and microSD subsystem.\\
	
	%% SPI bus
	\rowcolor{oddcol}
	\textbf{ICSP SCK} & D13 & J~11 & pin~23 & GPIO11/SPI0 SCLK &
	\fixed. 33\,$\Omega$ series resistor on PCB. 
	$\to$ TFT ILI9341 SCK.\\
	
	%% SPI bus
	\rowcolor{evencol}
	\textbf{TFT CS} & D10 & J~6 & pin~24 & GPIO8/SPI0 CE0 &
	\fixed\ 33\,$\Omega$ series resistor on PCB. Active-low. 
	$\to$ TFT ILI9341 CS --- Chip Select\\

	%% SPI bus
	\rowcolor{oddcol}
	\textbf{TFT DC} & D9 & J~7 & pin~31 & GPIO6 &
	\soft. 33\,$\Omega$ series resistor on PCB. 
	$\to$ TFT ILI9341 DC --- Data/Command select\\

	%% SD card
	\rowcolor{evencol}
	\textbf{SD CS} & D4 & J~8 & Pin~29 & GPIO5 &
	\soft. Active-low. 33\,$\Omega$ series resistor on PCB.
	SDcard Chip Select. (May change if microSD is not used.)\\

	%% Touch / I2C
	\rowcolor{oddcol}
	\textbf{SDA} & SDA(17) & J~2 & Pin~3 & GPIO2/I2C SDA1 &
	\fixed. 33\,$\Omega$ series resistor on PCB for ESD.  
	The 1651 shield has onboard 4.7\,k$\Omega$ pull-up.
	Do \emph{not} add a second pull-up on PCB.
	$\to$ TSC2007-SDA \\
	
	\rowcolor{evencol}
	\textbf{SCL} & SCL(18) & J~4 & Pin~5 & GPIO3/I2C SCL1 &
	\fixed. Same pull-up and ESD notes as SDA. 
	$\to$ TSC2007-SCL\\
	
	\rowcolor{oddcol}
	\textbf{TSIRQ} & D~2 & J~10 & Pin~7 & GPIO4 &
	\soft. Touch interrupt (active-low). Solder the TSIRQ/\#2 jumper on 1651 to enable.
	Enable RPi internal pull-up in software.
	33\,$\Omega$ series resistor on PCB. \\

	%% Backlight
	\rowcolor{evencol}
	\textbf{LITE} & D3 & J~9 &Pin~33 & GPIO13 &
	\soft. Backlight PWM. 
	Hardware PWM channel 1. 
	Solder the LITE/\#3 jumper on 1651 to enable.
	33\,$\Omega$ series resistor on PCB. \\
	
	%% VddIO note
	\rowcolor{warn}
	\textbf{VddIO jumper} &
	\multicolumn{5}{L{14cm}}{%
		\textbf{IMPORTANT:} The 1651 VddIO jumper must be set to \textbf{3V}
		(solder left pad to centre pad) because the RPi Zero 2\,W is a 3.3\,V logic device.
		This selects the logic reference voltage for the 74LVX245 level shifters on the shield.
		Incorrect setting risks damaged GPIO pins.} \\
	
\end{longtable}
% =============================================================
\section{Passive Component Recommendations for DIY Kit}
\label{sec:passive}

\subsection{Protective Series Resistors on 1651 Signal Lines}

All signal lines between the RPi GPIO header and the 1651 shield header
should have a \textbf{33\,$\Omega$ series resistor} placed on the PCB,
on the RPi side of the wire run (i.e., close to the RPi header).

\begin{notebox}
\textbf{Rationale:}
The RPi's GPIO pins have no internal current limiting or ESD protection diodes
rated for external misconnection.
A 33\,$\Omega$ resistor limits the peak fault current if a signal line is
accidentally shorted to 5\,V or GND during kit assembly or use,
and damps transmission-line ringing on the SPI clock and data lines
at high speed (40\,MHz).
At 40\,MHz SPI, the RC time constant with typical PCB trace capacitance
($\sim$5\,pF) gives a corner frequency well above 100\,MHz,
so the resistor has negligible effect on signal integrity at this speed.

\smallskip
\textbf{Value:} 33\,$\Omega$ is the standard recommended value for RPi GPIO protection.
Acceptable range: 22--47\,$\Omega$.
\textbf{Package:} 0402 or 0603 SMD preferred for a compact PCB.
Place as close as possible to the RPi GPIO header pin.
\end{notebox}

\textbf{Specific lines and resistor placement:}

\begin{itemize}[noitemsep]
  \item \textbf{SPI lines (MOSI, MISO, SCLK, TFT CS, TFT DC, SD CS):}
        33\,$\Omega$ each. Place within 5\,mm of RPi header.
  \item \textbf{I2C lines (SDA, SCL):}
        33\,$\Omega$ each. Note: the 1651 shield already has 4.7\,k$\Omega$ pull-ups
        on these lines. Do \emph{not} add PCB pull-ups --- the shield pull-ups are sufficient.
  \item \textbf{TSIRQ:} 33\,$\Omega$ series. RPi internal pull-up enabled in software.
  \item \textbf{LITE (backlight PWM):} 33\,$\Omega$ series.
  \item \textbf{5\,V power line:} No series resistor. Use adequate trace width
        ($\geq$0.5\,mm for up to 500\,mA).
  \item \textbf{GND:} No series resistor.
\end{itemize}

\subsection{1-Wire Bus Pull-up}

\begin{warnbox}
The 1-Wire bus on GPIO7 (physical pin~26) \textbf{requires a 4.7\,k$\Omega$ pull-up
resistor to 3.3\,V} on the PCB.
This is a hard requirement of the 1-Wire protocol --- the internal RPi pull-up
(50\,k$\Omega$) is far too weak and will cause communication failures.
Use a 4.7\,k$\Omega$ resistor to the 3.3\,V rail (RPi pin~1 or pin~17).
Package: 0402 or 0603 SMD.
\end{warnbox}

\subsection{LED Current-Limit Resistors}

All LED outputs (indicator LEDs on GPIO12, GPIO18, GPIO23;
control LEDs on GPIO16, GPIO20, GPIO21) must have current-limiting
resistors in series with the LED on the PCB.

\begin{itemize}[noitemsep]
  \item RPi GPIO output voltage: 3.3\,V
  \item Maximum current per GPIO pin: 16\,mA (recommended: 8--12\,mA for longevity)
  \item Typical LED forward voltage: 2.0\,V (red/yellow), 3.0\,V (blue/green/white)
  \item \textbf{Recommended resistor for red/yellow LED:}
        $(3.3 - 2.0)\,\text{V} / 10\,\text{mA} = 130\,\Omega$, use \textbf{150\,$\Omega$}
  \item \textbf{Recommended resistor for blue/green/white LED:}
        $(3.3 - 3.0)\,\text{V} / 10\,\text{mA} = 30\,\Omega$, use \textbf{33\,$\Omega$}
        (or higher if the LED is very bright)
  \item Package: 0402 or 0603 SMD.
\end{itemize}

\subsection{MOSFET Gate Resistors}

For MOSFET gate drive lines (GPIO16, GPIO20, GPIO21, and any remaining
MOSFET control pins), place a \textbf{33--100\,$\Omega$ gate resistor}
in series between the GPIO output and the MOSFET gate.
This limits the gate charge current spike at switch-on,
which would otherwise present a momentary short-circuit load to the GPIO pin.
If MOSFETs switch inductive loads, add a flyback diode across the load.

\subsection{Decoupling Capacitors}

\begin{notebox}
\textbf{On the RPi header power pins:}
Place a \textbf{100\,nF ceramic capacitor} (0402 or 0603, X5R or X7R)
between the 5\,V supply pin (pin~2) and the nearest GND pin (pin~6)
on the PCB, as close as possible to the connector.
This decouples switching transients on the supply caused by the display's
SPI bursts.

\smallskip
\textbf{On the 1651 shield header power pin:}
Place a \textbf{100\,nF ceramic} plus optionally a \textbf{10\,$\mu$F electrolytic}
(or tantalum) in parallel, between VIN and GND at the shield power header pins.
The larger capacitor absorbs the inrush current when the display backlight
and SPI bus first activate.

\smallskip
\textbf{On the 3.3\,V rail (for 1-Wire and I2C pull-ups):}
Place a \textbf{100\,nF ceramic} between 3.3\,V (pin~1 or 17) and GND
close to the pull-up resistors.

\smallskip
\textbf{No capacitors are needed} in the GPIO signal paths.
The 33\,$\Omega$ series resistors already damp any ringing.
Adding capacitance to SPI lines would slow the edges at high speed.
\end{notebox}

% =============================================================
\section{Software Configuration Notes}

\begin{itemize}[noitemsep]
  \item \textbf{SPI:} Enable with \texttt{dtparam=spi=on} in \texttt{/boot/config.txt}
  \item \textbf{I2C:} Enable with \texttt{dtparam=i2c\_arm=on} in \texttt{/boot/config.txt}
  \item \textbf{1-Wire:} Enable with
        \texttt{dtoverlay=w1-gpio,gpiopin=7} in \texttt{/boot/config.txt}
        (non-default pin; standard default is GPIO4)
  \item \textbf{Backlight PWM:} GPIO13 is hardware PWM channel~1 (PWM1).
        Use \texttt{pigpio} library for precise frequency and duty-cycle control,
        or \texttt{RPi.GPIO} for simple on/off and software PWM.
  \item \textbf{TSIRQ:} Configure GPIO4 as input with pull-up.
        TSC2007 default I2C address: \texttt{0x48}.
  \item \textbf{1651 VddIO jumper:} Must be set to 3V before powering on.
\end{itemize}

% =============================================================
\section{Summary: Unallocated and Reserved Pins}

\begin{tabular}{L{2cm} L{2cm} L{9cm}}
  \toprule
  \textbf{Physical} & \textbf{BCM} & \textbf{Status} \\
  \midrule
  Pin 1  & \na     & 3.3\,V --- available for low-current peripherals \\
  Pin 17 & \na     & 3.3\,V --- available for low-current peripherals \\
  Pin 13 & GPIO27  & \textit{Spare odd-side GPIO --- reserve} \\
  Pin 15 & GPIO22  & \textit{Spare odd-side GPIO --- reserve} \\
  Pin 27 & GPIO0   & EEPROM --- do not use \\
  Pin 28 & GPIO1   & EEPROM --- do not use \\
  Pin 8  & GPIO14  & UART TX --- avoid in kit \\
  Pin 10 & GPIO15  & UART RX --- avoid in kit \\
  \bottomrule
\end{tabular}

\vspace{10pt}
\begin{warnbox}
\textbf{Outstanding allocation items --- finalise before PCB layout:}
\begin{itemize}[noitemsep]
  \item 2 remaining MOSFET control outputs (of 5 total; 3 assigned: GPIO16, 20, 21)
  \item 2 remaining digital inputs (of 4 total; 2 assigned: GPIO24, 25)
  \item Candidates: even-side GPIO22 (pin~15 odd --- reconsidered if moved),
        or confirm further even-side spares with designer
  \item Confirm whether all 5 control LEDs vs MOSFET gate signals share the same pins,
        or whether LED indicators and MOSFET gates are on separate GPIO lines
\end{itemize}
\end{warnbox}

\end{document}
